\documentclass[shortabstract, inz]{iithesis}

\polishtitle    {Studium implementacji obsługi instrukcji \fmlinebreak force/release w Verilatorze}
\englishtitle   {A study of implementing handling of force/release statements in Verilator}

\polishabstract {Instrukcje \texttt{force} i \texttt{release} zdefiniowane w standardzie IEEE 1364-2001~\cite{IEEE1364-2001} języka Verilog umożliwiają wymuszenie wartości sygnałów w symulacji. W pracy przedstawiono różne implementacje obsługi tych instrukcji w otwartoźródłowym symulatorze Verilator, w tym opracowane przez autora niniejszej pracy. Opisano semantykę działania instrukcji \texttt{force}/\texttt{release}, a także omówiono szczegóły implementacji oraz testy poprawności działania.}
\englishabstract{The \texttt{force} and \texttt{release} statements defined in the IEEE 1364-2001~\cite{IEEE1364-2001} Verilog standard allow forcing signal values in simulation. This thesis presents various implementations of these statements in the open-source Verilator simulator, including those developed by the author of this thesis. The semantics of the \texttt{force}/\texttt{release} statements are described, along with details of the implementation and correctness tests.}
\author         {Artur Bieniek}
\advisor        {mgr inż. Kuba Nowak}
\date          {5 lutego 2026}                     % Data zlozenia pracy
% Dane do oswiadczenia o autorskim wykonaniu
\transcriptnum {347165}                     % Numer indeksu
\advisorgen    {mgr. Kuby Nowaka} % Nazwisko promotora w dopelniaczu
%%%%%

%%%%% WLASNE DODATKOWE PAKIETY
\usepackage[utf8]{inputenc}
\usepackage{xcolor}
\usepackage{minted}
\usepackage{tabularx}
\usepackage{booktabs}
\usepackage{float}
\usepackage{caption}
\usepackage{fvextra}

% Czcionka taka żeby tyldy były na środku
\usepackage[T1]{fontenc}
\usepackage{lmodern}


\fvset{breaklines=true, breakanywhere=true}
% Nie zostawiaj samotnej linii akapitu na stronie
\widowpenalty=10000
\clubpenalty=10000
\displaywidowpenalty=10000

% --- Tabele z auto-liczeniem zmian [%] i znakiem + / bez znaku dla 0 ---

\usepackage{xfp}      % obliczenia zmiennoprzecinkowe: \fpeval{...}
\usepackage{siunitx}  % formatowanie liczb: \num{...}

\sisetup{
  round-mode      = places,
  round-precision = 2,
  round-pad       = true,
  output-decimal-marker = {.}
}

% (impl - ref) / ref * 100, zaokrąglone do 2 miejsc
% dodatnie: +X.XX, ujemne: -X.XX, zero: 0.00 (bez znaku)
\newcommand{\pctchangeSigned}[2]{%
  \ifnum\fpeval{round(((#2-#1)/#1)*100, 2) > 0}=1
    +\num{\fpeval{round(((#2-#1)/#1)*100, 2)}}%
  \else
    \ifnum\fpeval{round(((#2-#1)/#1)*100, 2) < 0}=1
      \num{\fpeval{round(((#2-#1)/#1)*100, 2)}}%
    \else
      \num{0.00}%
    \fi
  \fi
}

% --- Podpisy "Kod X.Y" bez caption/listing (nie psuje layoutu) ---
\newcounter{kod}[chapter]
\renewcommand{\thekod}{\thechapter\arabic{kod}}

\newcommand{\codecaption}[2]{%
  \par\refstepcounter{kod}\label{#1}%
  \noindent\textbf{Kod~\thekod:}~#2\par
  \vspace*{1.0\baselineskip}%
}

\newcommand{\coderef}[1]{\hyperref[#1]{Kod~\ref*{#1}}}

% Wiersz: label/design, Tref, Rref, Timpl, Rimpl
% Wypisuje dane + oblicza dwie zmiany [%] bez duplikowania liczb.
\newcommand{\benchrow}[5]{%
#1 & #2 & #3 & #4 & #5 & \pctchangeSigned{#2}{#4} & \pctchangeSigned{#3}{#5} \\ \hline
}

% Pozwala bardziej łamać linie, żeby rzeczy nie wychodziły na margines
\AtBeginDocument{\sloppy}

\usepackage[backend=biber,style=numeric,sorting=none]{biblatex}
\addbibresource{citations.bib}

\newcommand{\verilatorfootnote}{\textbf{Uwaga:} Ilekroć w pracy pojawia się stwierdzenie \textit{obecnie w Verilatorze}, odnosi się ono do stanu implementacji na dzień 2 listopada 2025 roku, wersji Verilator 5.042.}

\makeatletter
\def\ps@verilatornote{%
  \let\@oddhead\@empty
  \let\@evenhead\@empty
  \def\@oddfoot{\parbox{\textwidth}{\raggedright\footnotesize\verilatorfootnote}}%
  \let\@evenfoot\@oddfoot
}
\makeatother

\begin{document}

%%%%% POCZĄTEK ZASADNICZEGO TEKSTU PRACY

\chapter{Wprowadzenie}

Verilog i SystemVerilog są popularnymi językami opisu sprzętu (HDL -- \textit{Hardware Description Language}), szeroko stosowanymi w projektowaniu, weryfikacji oraz symulacji układów cyfrowych. Umożliwiają one opis projektu na poziomie RTL (\textit{Register Transfer Level}), gdzie zachowanie układu definiowane jest na poziomie przepływu danych pomiędzy rejestrami. Następnie narzędzia syntezy przekształcają taki opis do postaci struktury złożonej z przerzutników, zatrzasków i bramek logicznych. Istotną zaletą podejścia RTL jest możliwość szybkiej weryfikacji projektu poprzez symulację komputerową na wczesnym etapie prac.

\section{Instrukcje \texttt{force}/\texttt{release} w weryfikacji układów}

Standard IEEE~1364-2001~\cite{IEEE1364-2001} (Verilog) w sekcji 10.6 (\textit{Procedural continuous assignments}) definiuje instrukcje \texttt{force} i \texttt{release} oraz \texttt{assign} i \texttt{deassign}. Mechanizmy te umożliwiają realizację proceduralnego przypisania o charakterze kombinacyjnym. W odróżnieniu od zwykłego przypisania kombinacyjnego (\texttt{assign}), przypisanie proceduralne jest aktywowane w kontekście procedury (np. w bloku \texttt{initial} lub \texttt{always}) i obowiązuje od wykonania instrukcji \texttt{force} (lub \texttt{assign}) aż do wykonania \texttt{release} (lub \texttt{deassign}).

Instrukcje \texttt{force}/\texttt{release} są szczególnie użyteczne w testach, gdzie służą do czasowego wymuszania określonych wartości sygnałów podczas symulacji. Pozwala to testować scenariusze trudne do uzyskania wyłącznie poprzez sterowanie portami wejściowymi, w tym również wymuszać wartości sygnałów wewnętrznych w hierarchii projektu.

\section{Verilator jako kontekst i motywacja pracy}

W klasycznych symulatorach sterowanych zdarzeniami każda zmiana wartości sygnału powoduje powstanie zdarzenia aktualizacji (ang. \textit{update event}), które może wyzwalać ponowną ewaluację procesów wrażliwych na ten sygnał. W takim modelu mechanizmy \texttt{force} i \texttt{release} są naturalnie powiązane z harmonogramowaniem zdarzeń i rozstrzyganiem priorytetów pomiędzy różnymi źródłami sterowania sygnałem.

Verilator~\cite{Verilator} jest natomiast otwartoźródłowym symulatorem opartym na kompilacji opisu sprzętu do kodu języka C++, który następnie jest kompilowany i wykonywany jako program natywny. Podejście to zapewnia wysoką wydajność symulacji, jednak nie opiera się bezpośrednio na klasycznym, zdarzeniowym kernelu symulacyjnym. W konsekwencji implementacja instrukcji \texttt{force}/\texttt{release} w Verilatorze wymaga zaprojektowania dodatkowych mechanizmów, które zapewnią zgodność obserwowalnego zachowania z semantyką języka.

Ze względu na fakt, że Verilator jest obecnie najszybszym i najczęściej wykorzystywanym otwartoźródłowym symulatorem RTL, istotne jest, aby wspierał możliwie szeroki zakres konstrukcji języka Verilog i SystemVerilog. Umożliwia to większej liczbie użytkowników korzystanie z zaawansowanych mechanizmów weryfikacyjnych w środowisku open-source. W aktualnych wersjach Verilatora obsługa instrukcji \texttt{force}/\texttt{release} jest jednak istotnie ograniczona, co stanowiło bezpośrednią motywację do podjęcia prac nad jej rozszerzeniem.

\section{Cel i zakres pracy}

Celem niniejszej pracy jest zbadanie możliwości implementacji instrukcji \texttt{force} i \texttt{release} w otwartoźródłowym symulatorze Verilator. Praca obejmuje analizę wymagań wynikających ze standardu opisującego język Verilog/SystemVerilog, projekt mechanizmu obsługi tych instrukcji, jego implementację oraz weryfikację poprawności i wpływu na wydajność symulacji. Przedstawione zostaną alternatywne podejścia implementacyjne, omówione zostaną napotkane problemy oraz zaprezentowane wyniki testów wraz z wnioskami dotyczącymi użyteczności i efektywności zaproponowanych rozwiązań.

\section{Przyjęta semantyka i odniesienie do modelu IEEE}

Choć instrukcje \texttt{force}/\texttt{release} zostały wprowadzone w Verilogu już w standardzie IEEE~1364-2001~\cite{IEEE1364-2001}, standard SystemVerilog (IEEE~1800-2023)~\cite{IEEE1800-2023} również je zawiera i definiuje ich semantykę. Ponieważ IEEE~1800-2023 jest najnowszym standardem języka, a Verilator docelowo dąży do możliwie pełnego wsparcia SystemVerilog, w niniejszej pracy przyjęto semantykę instrukcji \texttt{force}/\texttt{release} zgodną ze standardem IEEE~1800-2023~\cite{IEEE1800-2023}.

Specyfikacja IEEE~1800-2023~\cite{IEEE1800-2023} opisuje semantykę SystemVerilog w oparciu o model wykonania bazujący na dyskretnych zdarzeniach oraz przedstawia referencyjny, abstrakcyjny algorytm symulacji. Algorytm ten określa porządek przetwarzania zdarzeń oraz zależności czasowe pomiędzy regionami symulacyjnymi. Standard dopuszcza jednak dowolność w doborze algorytmów i struktur danych w implementacji symulatora, pod warunkiem że jego obserwowalne zachowanie pozostaje zgodne z zachowaniem wynikającym z algorytmu referencyjnego.

W związku z tym prezentowane w pracy implementacje instrukcji \texttt{force}/\texttt{release} w Verilatorze mają na celu odwzorowanie zachowania zgodnego z modelem zdarzeniowym IEEE, pomimo że sam Verilator nie realizuje tego modelu bezpośrednio w postaci klasycznego, zdarzeniowego kernela symulacyjnego.

\section{Geneza pracy i kontekst realizacji}
Część inżynierska niniejszej pracy została zrealizowana w ramach pracy zawodowej autora w firmie Antmicro, zajmującej się projektowaniem i weryfikacją systemów elektronicznych oraz rozwojem narzędzi open-source. Prace badawczo-inżynierskie prowadzone w tym kontekście obejmowały rozwój implementacji obsługi instrukcji force/release w symulatorze Verilator. Jednocześnie całość tekstu pracy dyplomowej, w tym dobór zakresu, analiza problemu, opis rozwiązań oraz wnioski, stanowi samodzielne opracowanie autora.

\thispagestyle{verilatornote}

\chapter{Wprowadzenie do semantyki przypisań}

\section{Wstęp}
Języki opisu sprzętu, takie jak Verilog i SystemVerilog, oferują różnorodne mechanizmy przypisywania wartości sygnałom, które różnią się semantyką oraz zastosowaniem. Zrozumienie tych mechanizmów jest kluczowe zarówno dla prawidłowego modelowania układów cyfrowych, jak i dla implementacji symulatorów wspierających te konstrukcje językowe. W niniejszym rozdziale przedstawiono systematyczny przegląd typów przypisań występujących w standardzie SystemVerilog, ze szczególnym uwzględnieniem ich zachowania w kontekście symulacji. Omówione zostaną kolejno:
\begin{itemize}
\item przypisania kombinacyjne, które modelują syntezowalne bezpośrednie połączenia sygnałów,
\item przypisania proceduralne, wykonywane jednokrotnie w ramach bloków proceduralnych --- syntezowalne jako przerzutniki, zatrzaski, układy kombinacyjne itp.,
\item przypisania proceduralne kombinacyjne realizowane za pomocą instrukcji \linebreak \texttt{assign}/\texttt{deassign} i \texttt{force}/\texttt{release} --- niesyntezowalne.
\end{itemize}

Te ostatnie stanowią szczególnie istotny element testowania układów cyfrowych, umożliwiając wymuszanie wartości sygnałów w dowolnym momencie symulacji. Zrozumienie hierarchii priorytetów między różnymi typami przypisań oraz ich wzajemnych interakcji jest niezbędne do prawidłowej implementacji mechanizmu \texttt{force}/\texttt{release} w symulatorze Verilator.

\section{Przypisanie kombinacyjne}
W kontekście sprzętu, przypisanie kombinacyjne opisuje syntezowalne bezpośrednie połączenie obu sygnałów. W kontekście symulacji, powinno zachowywać się identycznie --- lewa strona przypisania powinna wskazywać przy każdym odczycie wartość identyczną jak prawa. Przykładowo, następujące przypisanie kombinacyjne:
\begin{minted}{systemverilog}
module top (input a, input b, output out);
    assign out = a & b;
endmodule
\end{minted}
\codecaption{lst:comb-assign}{Przypisanie kombinacyjne w Verilogu/SystemVerilogu.}
spowoduje, że sygnał \texttt{out} będzie przy każdym odczycie równy iloczynowi logicznemu sygnałów \texttt{a} i \texttt{b}. Gdy tylko wartość \texttt{a} lub \texttt{b} się zmieni, wartość \texttt{out} zostanie automatycznie zaktualizowana.

Innym sposobem przypisania kombinacyjnego jest deklaracja sieci~\footnote{sieci (ang. \textit{net}), deklarowane za pomocą słów kluczowych \texttt{wire}, \texttt{tri}, \texttt{wand}, \texttt{wor} itp., służą do modelowania połączeń fizycznych między elementami układu cyfrowego. W odróżnieniu od zmiennych, które przechowują wartości i mogą być przypisywane proceduralnie, sieci reprezentują sygnały, których wartość jest określana przez przypisania kombinacyjne.}

\begin{minted}{systemverilog}
module top (input a, input b);
    wire out = a & b; // deklaracja sieci z inicjalizacją
    initial $monitor("%b", out);
endmodule
\end{minted}
\codecaption{lst:net-decl-assign}{Deklaracja sieci z inicjalizacją w Verilogu/SystemVerilogu.}
W tym przykładzie, sygnał \texttt{out} jest przypisany za pomocą przypisania kombinacyjnego do iloczynu logicznego sygnałów \texttt{a} i \texttt{b}. Wartości obserwowane na monitorze w kolejnych krokach symulacji będą odzwierciedlać iloczyn logiczny wartości \texttt{a} i \texttt{b}.

\section{Przypisanie proceduralne}
Standardowe przypisanie proceduralne w Verilogu/SystemVerilogu, takie jak:
\begin{minted}{systemverilog}
module top (input a, input b, output reg out);
    always @(*) begin
        out = a;
        out = out & b;
    end
endmodule
\end{minted}
\codecaption{lst:proc-assign}{Przypisanie proceduralne w Verilogu/SystemVerilogu.}

jest analizowane przez narzędzie do syntezy, które na tej podstawie wnioskuje:
\begin{itemize}
    \item układy kombinacyjne gdy 
    \begin{itemize}
        \item każde wyjście jest przypisane dla wszystkich możliwych przypadków
        \item wartość wyjścia zależy tylko i wyłącznie od aktualnych wartości wejść
        \item nie ma żadnego wewnętrznego stanu przechowywanego między kolejnymi ewaluacjami
    \end{itemize}
    \item zatrzaski gdy
    \begin{itemize}
        \item sygnał nie jest przypisany w każdej możliwej ścieżce wykonania
        \item blok, w którym występuje przypisanie, nie jest wyzwalany przez sygnały zegarowe
    \end{itemize}
    \item przerzutniki gdy
    \begin{itemize}
        \item blok, w którym występuje przypisanie, jest wyzwalany przez sygnały zegarowe
        \item wartość sygnału zmienia się tylko w odpowiedzi na zbocza sygnału zegarowego
        \item pomiędzy zboczami sygnału zegarowego stan jest przechowywany
    \end{itemize}
    \item multipleksery dla
    \begin{itemize}
        \item operatora ternarnego: \texttt{warunek ? wartosc\_a : wartosc\_b}
        \item instrukcji warunkowej: \texttt{wartosc\_a if warunek else wartosc\_b}
        \item instrukcji \texttt{case}
    \end{itemize}
    \item inne konstrukcje, dostępne w danym procesie produkcyjnym lub na danym układzie FPGA, takie jak liczniki, pamięci, układy arytmetyczne itp.
\end{itemize}
W symulacji przypisanie to jest wykonywane tylko raz, w momencie wykonania instrukcji (podczas wywoływania zdarzenia w przypadku symulatorów sterowanych zdarzeniami lub podczas przejścia przez tę instrukcję w skompilowanym kodzie w przypadku Verilatora). Wartość sygnału \texttt{out} nie jest automatycznie aktualizowana w odpowiedzi na zmiany wartości \texttt{a} lub \texttt{b}. Aby zaktualizować wartość \texttt{out}, instrukcja musi zostać ponownie wykonana, na przykład w ramach bloku \texttt{always} reagującego na zmiany sygnałów.

Przypisania proceduralne są szczególnie użyteczne w przypadku projektowania układów sekwencyjnych, gdzie stan układu jest przechowywany i aktualizowany w odpowiedzi na sygnały zegarowe lub inne zdarzenia. Pozwalają one na modelowanie złożonych zachowań, które nie mogą być łatwo wyrażone za pomocą prostych przypisań kombinacyjnych. Zachowanie układu można opisywać wówczas jak ciąg instrukcji wykonywanych sekwencyjnie --- choć w rzeczywistości tak nie jest, zachowanie jest identyczne.

Obecnie w Verilatorze przypisania proceduralne są realizowane poprzez bezpośrednie przypisania w kodzie C++ wygenerowanym z kodu Verilog/SystemVerilog. Oznacza to, że wartość sygnału jest aktualizowana wtedy, gdy odpowiednia instrukcja jest wykonywana w ramach pętli symulacyjnej.

\section{Przypisanie proceduralne kombinacyjne}
Przypisanie proceduralne kombinacyjne za pomocą instrukcji \texttt{assign} łączy cechy przypisań kombinacyjnych i proceduralnych. Instrukcja \texttt{assign} pozwala na wykonanie przypisania, które zachowuje się jak przypisanie kombinacyjne, ale jest wykonywane w kontekście procedury. Oznacza to, że przypisanie to jest aktywne od momentu jego wykonania do momentu wyłączenia go za pomocą instrukcji \texttt{deassign}.
Jeśli wartość jakiejkolwiek zmiennej z prawej strony takiego przypisania zmieni się, to powinna nastąpić ponowna ewaluacja, o ile przypisanie to jest aktywne (czyli zostało wykonane instrukcją \texttt{assign} i nie zostało jeszcze wyłączone instrukcją \texttt{deassign}).
Przykładowo:
\begin{minted}{systemverilog}
module top (input a, input b, output reg out);
    initial begin
        out = 0;
        assign out = a & b;
        out = 1; // nie ma efektu
        $display("out = %b", out);
    end
endmodule
\end{minted}
\codecaption{lst:proc-cont-assign}{Przypisanie proceduralne kombinacyjne w Verilogu/SystemVerilogu.}

Przypisanie za pomocą \texttt{assign} ma wyższy priorytet niż zwykłe przypisanie proceduralne. Oznacza to, że jeśli sygnał jest przypisany zarówno za pomocą \texttt{assign}, jak i zwykłego przypisania kombinacyjnego lub proceduralnego, to wartość z \texttt{assign} ma pierwszeństwo.
Dlatego przypisanie proceduralne następujące później nie ma żadnego efektu, ponieważ przypisanie za pomocą \texttt{assign} jest nadal aktywne i ma wyższy priorytet.
Jeśli wartość \texttt{a} lub \texttt{b} się zmieni, to wartość \texttt{out} powinna zostać ponownie obliczona.
Jeśli przypisanie proceduralne kombinacyjne zostanie wykonane do zmiennej, do której jest aktywne inne przypisanie proceduralne kombinacyjne, to nowe przypisanie powinno wyłączyć poprzednie, zanim stanie się aktywne.
Przypisanie proceduralne za pomocą \texttt{assign} może być wykonane tylko do pojedynczej zmiennej lub do konkatenacji zmiennych, w szczególności nie może być wykonane do elementu tablicy ani do części zmiennej (np. wycinka wektora bitów).

Użycie słowa kluczowego \texttt{assign} w dwóch różnych kontekstach (przypisanie kombinacyjne oraz proceduralne kombinacyjne) może powodować dezorientację --- więcej o tym w \ref{sec:deassign-deprecation}

\coderef{lst:proc-cont-assign} pokazuje przypadek, którego Verilator obecnie nie obsługuje poprawnie --- przypisania proceduralne kombinacyjne są obecnie traktowane jak zwykłe przypisania proceduralne, bez automatycznej ponownej ewaluacji przy zmianie wartości zmiennych z prawej strony przypisania. W efekcie, wartość \texttt{out} pozostanie równa 1, niezależnie od wartości \texttt{a} i \texttt{b}.

\subsection{Instrukcja \texttt{deassign}}
Instrukcja \texttt{deassign} służy do wyłączenia aktywnego przypisania proceduralnego kombinacyjnego wykonanego wcześniej za pomocą instrukcji \texttt{assign}. Po wykonaniu instrukcji \texttt{deassign}, przypisanie przestaje być aktywne, a tym samym wartość sygnału przestaje być aktualizowana. Wartość sygnału pozostaje bez zmian do momentu przypisania mu nowej wartości za pomocą przypisania proceduralnego lub proceduralnego kombinacyjnego.
Ta instrukcja nie jest aktualnie obsługiwana w Verilatorze, jej użycie skutkuje błędem.
\begin{verbatim}
%Error-UNSUPPORTED: test.v:11:12: Unsupported: Verilog 1995 deassign
11 |         #1 deassign r;
   |            ^~~~~~~~
\end{verbatim}
\codecaption{lst:deassign-error}{Błąd generowany przez Verilator przy użyciu instrukcji \texttt{deassign}.}


Zarówno instrukcja \texttt{assign} wykonana w kontekście procedury, innej niż blok \texttt{always}, jak i instrukcja \texttt{deassign} nie są syntezowalne, ponieważ nie posiadają odpowiednika w sprzęcie i służą wyłącznie do oddziaływania na zachowanie testowanego układu w symulacji.

\section{Przypisanie za pomocą \texttt{force}}
Instrukcja \texttt{force} ma podobny efekt do instrukcji \texttt{assign}, ale może być użyta do wymuszenia wartości również sieci, a nie tylko zmiennych. Ponadto lewa strona przypisania za pomocą \texttt{force} może być bardziej złożona, na przykład może zawierać elementy tablicy lub wycinki wektorów bitów (o stałych granicach). Nie może jednak być użyte do przypisania wartości do zmiennej która została przypisana za pomocą mieszanki przypisań kombinacyjnych i proceduralnych.
Przypisanie \texttt{force} do zmiennej, dopóki jest aktywne, przykrywa przypisania proceduralne, kombinacyjne oraz proceduralne kombinacyjne.
W oczywisty sposób, przypisanie \texttt{force} nie jest syntezowalne, ponieważ nie posiada odpowiednika w sprzęcie i służy wyłącznie do sterowania przebiegiem symulacji.

\noindent
\begin{tabular}{@{}p{0.6\linewidth}@{\hspace{1em}}p{0.34\linewidth}@{}}
\textbf{Kod} & \textbf{Wynik symulacji}\\ \hline

\begin{minipage}[t]{\linewidth}
\begin{minted}[linenos,breaklines,fontsize=\small]{systemverilog}
module top;
    reg a=0,b=1;
    reg [1:0] r;

    initial begin
        $monitor("%0t a=%b, b=%b, r=%b", $time, a, b, r);
        r = 1'b0;
        #1 assign r = 1'b1;
        #1 r = 1'b0; // ignorowane
        #1 deassign r;
        #1 r = 1'b0;
        #1 assign r = a;
        #1 a = 1'b1;
        #1 a = 1'b0;
        #1 force r = a + b;
        #1 a = 1'b0; b = 1'b0;
        #1 a = 1'b1; b = 1'b1;
        #1 assign r = b; // przykryte
        #1 r = 2'b11; // ignorowane
        #1 release r;
        #1 b = 1'b0;
        #1 $finish;
    end
endmodule
\end{minted}
\end{minipage}
&
\begin{minipage}[t]{\linewidth}
\begin{minted}[fontsize=\small]{text}







0 a=0, b=1, r=00
1 a=0, b=1, r=01


4 a=0, b=1, r=00

6 a=1, b=1, r=01
7 a=0, b=1, r=00
8 a=0, b=1, r=01
9 a=0, b=0, r=00
10 a=1, b=1, r=10


13 a=1, b=1, r=01
14 a=1, b=0, r=00
\end{minted}
\end{minipage}
\\
\end{tabular}
\codecaption{lst:hierarchy}{Zilustrowanie hierarchii priorytetów przypisań w Verilogu/SystemVerilogu. Instrukcja w linii 7 to zwykłe przypisanie proceduralne, które się wykonuje. Następnie w linii 8 występuje przypisanie proceduralne kombinacyjne, które przykrywa poprzednie przypisanie. Kolejne przypisanie proceduralne w linii 9 jest ignorowane, ponieważ przypisanie proceduralne kombinacyjne jest nadal aktywne. Po wywołaniu \texttt{deassign} w linii 10, przypisanie proceduralne kombinacyjne zostaje wyłączone i kolejne przypisanie proceduralne w linii 11 ma efekt. W linii 12 ponownie wykonywane jest przypisanie proceduralne kombinacyjne. W liniach 13-14 zmieniana jest wartość prawej strony przypisania \texttt{assign}, co powoduje aktualizację wartości \texttt{r}. W linii 15 wykonywane jest przypisanie za pomocą \texttt{force}, które przykrywa wszystkie poprzednie przypisania. W liniach 16-17 zmieniane są wartości \texttt{a} i \texttt{b}, co powoduje aktualizację wartości \texttt{r} zgodnie z przypisaniem \texttt{force}. W linii 18 wykonywane jest przypisanie proceduralne kombinacyjne, które jest przykryte, ponieważ przypisanie \texttt{force} jest nadal aktywne. Przypisanie proceduralne z linii 19 również jest przykryte. Po wywołaniu \texttt{release} w linii 20, przypisanie \texttt{force} zostaje wyłączone, na skutek czego przypisanie proceduralne kombinacyjne z linii 18 zostaje odkryte i reaguje na zmianę wartości \texttt{b} w linii 21.}

\subsection{Wycofanie instrukcji \texttt{assign}/\texttt{deassign}}
\label{sec:deassign-deprecation}
\medskip
Jak można zauważyć, para instrukcji \texttt{force}/\texttt{release} rozszerza funkcjonalność instrukcji \texttt{assign}/\texttt{deassign}, umożliwiając wymuszanie wartości również na sieciach oraz bardziej złożonych lewych stronach przypisań. W szczególności każde przypisanie za pomocą \texttt{assign}/\texttt{deassign} można wyrazić poprzez \texttt{force}/\texttt{release}. Ponadto, użycie słowa kluczowego \texttt{assign} w dwóch różnych kontekstach (przypisanie kombinacyjne oraz proceduralne kombinacyjne) może powodować dezorientację. Z tego powodu, już począwszy od IEEE 1800-2005~\cite{IEEE1800-2005} zaznaczono, iż instrukcje te mogą wkrótce zostać wycofane i kolejne wersje standardu mogą ich już nie wymagać. W IEEE 1800-2023~\cite{IEEE1800-2023} nadal wymaga się obsługi tych instrukcji, jednak stanowczo odradza się ich użycia na rzecz \texttt{force}/\texttt{release} i wciąż rozważane jest ich wycofanie.



\chapter{Opis implementacji w Verilatorze}
\section{Obecny stan}
Przed rozpoczęciem implementacji w ramach tej pracy, Verilator nie obsługiwał instrukcji \texttt{deassign}. Obsługa instrukcji \texttt{force} i \texttt{release} była częściowo zaimplementowana, jednak miała pewne ograniczenia, wynikające ze sposobu działania tejże implementacji.

Za obsługę \texttt{force} i \texttt{release} w Verilatorze odpowiada plik \texttt{V3Force.cpp}. Implementacja ta polega na trzymaniu dla każdego sygnału do którego jest wykonywane jakieś przypisanie \texttt{force} dodatkowych sygnałów:
\begin{itemize}
    \item \texttt{<v>\_\_VforceEn} --- zmienna typu logicznego, takiego samego typu jak sygnał lub pojedynczy bit w przypadku zmiennych skalarnych,~\footnote{Zmienne skalarne to takie, których nie można przypisywać częściowo używając wycinka bitowego, np. \textit{real}, \textit{int},  \textit{time}. Chociaż takie zmienne często w swojej reprezentacji składają się z wielu bitów, to takie przypisanie nie ma większego sensu i zgodnie ze specyfikacją nie jest możliwe} przechowująca informację o tym, czy dla danego bitu zmiennej jest aktywne przypisanie \texttt{force},
    \item \texttt{<v>\_\_VforceVal} --- zmienna tego samego typu jak sygnał, przechowująca ostatnią wymuszoną wartość,
    \item \texttt{<v>\_\_VforceRd} --- sieć takiego samego typu jak sygnał, aktualizowana na bieżąco do wartości \texttt{<v>\_\_VforceEn \& <v>\_\_VforceVal | ($\sim$<v>\_\_VforceEn) \& v} w przypadku zmiennych nieskalarnych lub do wartości \texttt{<v>\_\_VforceEn ? <v>\_\_VforceVal : v} w przypadku zmiennych skalarnych, przy czym \texttt{v} oznacza wartość sygnału bez wymuszenia

\end{itemize}
% Moim zdaniem łączenie tego z akapitem poprzedzającym enumerację nie jest dobre, ponieważ mówi o sygnale, który jeszcze nie został wprowadzony
Wszystkie wystąpienia sygnału, do którego jest wykonywane przypisanie \texttt{force}, są zastępowane przez \texttt{\_\_VforceRd}.

Dodatkowo, po każdym przypisaniu proceduralnym do którejkolwiek ze zmiennych z prawej strony przypisania \texttt{force}, wykonywana jest aktualizacja sygnału \texttt{\_\_VforceVal}, aby uwzględnić zmiany wartości tych zmiennych, co zapewnia poprawność późniejszego odczytu.

To podejście ma oczywiste ograniczenie: interpretacja statyczna kodu polega jedynie na najpóźniej występującej instrukcji \texttt{force}, więc nie obsługuje poprawnie przypadków, gdy instrukcja \texttt{force} jest wykonana pod wyrażeniem warunkowym, co może prowadzić do niepoprawnych wyników symulacji.

\begin{minted}{systemverilog}
`define IMPURE_ONE |($random | $random)
module top;
  wire a;
  reg b;
  initial begin
    if (`IMPURE_ONE == 1) force a = 1;
    if (`IMPURE_ONE == 0) force a = b;
    b = 0;
    #1 $display(a);
    release a;
  end
endmodule
\end{minted}
\codecaption{lst:conditional-force}{Przykład kodu, gdzie wymagana jest obsługa instrukcji \texttt{force} pod wyrażeniem warunkowym. Aktualizacja wartości sygnału \texttt{b} nie powinna mieć wpływu na wartość sygnału \texttt{a}, ponieważ przypisanie \texttt{force a = b} nie zostało wykonane}

\coderef{lst:conditional-force} zostanie przekształcony na:
\begin{minted}[linenos]{systemverilog}
logic [1:0] a;
logic [1:0] a__VforceRd;
logic [1:0] a__VforceEn = 2'h0;
logic [1:0] a__VforceVal;
logic [1:0] b;
always @(a__VforceRd or a__VforceVal or a__VforceEn or a)
    a__VforceRd = (a__VforceEn & a__VforceVal) | (~a__VforceEn & a);
initial begin
    if (`IMPURE_ONE) begin
        a__VforceEn = 2'h3;
        a__VforceVal = 2'h1;
        a__VforceRd = 2'h1;
    end

    if (~`IMPURE_ONE) begin
        a__VforceEn = 2'h3;
        a__VforceVal = b;
        a__VforceRd = b;
    end

    b = 2'h0;
    a__VforceVal = b;

    #1 $display("%d", a__VforceRd);

    a = (a__VforceEn & a__VforceVal) | ((~a__VforceEn) & a);
    a__VforceRd = (a__VforceEn & a__VforceVal) | ((~a__VforceEn) & a);
    a__VforceEn = 2'h0;
end
\end{minted}
\codecaption{lst:force-current}{Przekształcenie wykonane przez Verilator, pozyskane przy użyciu flag \texttt{----dump-tree ----debug-emitv}, uproszczone i sformatyzowane ręcznie dla czytelności. Linie 1-7 definiują i inicjalizują dodatkowe sygnały oraz zawierają blok \texttt{always} odpowiedzialny za aktualizowanie sygnału \texttt{\_\_VforceRd}. Linie 10-12 odpowiadają pierwszemu przypisaniu \texttt{force}, linie 16-18 drugiemu, linie 21-22 aktualizują wartość \texttt{\_\_VforceVal} w odpowiedzi na zmianę sygnału \texttt{b}, który był prawą stroną ostatniego napotkanego przypisanie \texttt{force}, a linie 26-28 realizują instrukcję \texttt{release}.}

Konstrukcja \texttt{IMPURE\_ONE}, wykorzystywana w wielu testach Verilatora, sama w sobie jest warta uwagi --- tworzy ona wyrażenie, które ma bardzo wysokie prawdopodobieństwo być równe 1, ale nie jest możliwe statyczne określenie jego wartości przed symulacją, więc symulator nie jest w stanie go wyoptymalizować. Wyrażenie to bierze dwie kolejne liczby losowe wygenerowane przez PRNG i wykonuje na nich operację OR bitową, a następnie wykonuje operację OR-redukcji na wyniku, co daje wynik 0 wtedy i tylko wtedy, gdy obie liczby losowe są równe 0, co ma prawdopodobieństwo $\frac{1}{2^{64}}$.

Wyżej opisana implementacja, choć niepoprawna w niektórych przypadkach, jest bardzo wydajna --- wymaga jedynie dodatkowych trzech sygnałów na każdy sygnał z wymuszeniem, a sama aktualizacja wartości wymuszonej jest wykonywana tylko wtedy, gdy któraś ze zmiennych z prawej strony przypisania ulegnie zmianie.

Przed rozpoczęciem pracy nad innymi implementacjami, dokonałem również kilku poprawek w tej obecnej, co powiększyło zakres poprawnie obsługiwanych przypadków oraz pomogło zagłębić się w jej działanie:
\begin{itemize}
    \item dodałem wsparcie dla wyrażenia z prawej strony \texttt{force} złożonego z wielu zmiennych \footnote{https://github.com/verilator/verilator/pull/6163},
    \item poprawiłem wyświetlanie błędów przy przypisaniu proceduralnym do wejścia, do którego istnieje przypisanie \texttt{force} \footnote{https://github.com/verilator/verilator/pull/6169},
    \item poprawiłem zachowanie symulacji, gdy jedna zmienna jest prawą stroną więcej niż jednego przypisania \texttt{force} \footnote{https://github.com/verilator/verilator/pull/6269},
    \item naprawiłem błąd \textit{Inconsistent assignment} generowany z powodu działania V3Force w wyniku niedopasowania szerokości lewej strony do prawej \footnote{https://github.com/verilator/verilator/pull/6541}.
\end{itemize}

Wprowadzone poprawki nie tylko zwiększyły praktyczną użyteczność istniejącej implementacji, lecz także unaoczniły jej ograniczenia wynikające z przyjętego modelu interpretacji statycznej, co bezpośrednio zmotywowało poszukiwanie alternatywnych podejść opisanych w kolejnych rozdziałach.

\section{Zastąpienie odczytów złożonymi wyrażeniami}
Pierwszą próbą implementacji~\footnote{https://github.com/ArturBieniek4/verilator/tree/81646-force-rework-v3} było zastąpienie każdego odczytu sygnału, do którego jest wykonywane przypisanie \texttt{force}, złożonym wyrażeniem, które uwzględnia wszystkie przypisania \texttt{force} do tego sygnału oraz to, czy są one aktualnie aktywne. W tym celu:
\begin{itemize}
    \item napotkane przypisania \texttt{force} do danego sygnału są rejestrowane w strukturze danych wraz z wyrażeniem stanowiącym ich prawą stronę oraz nadanym im unikalnym identyfikatorem \texttt{ID},
    \item dla każdego przypisania \texttt{force <v> = <RHS>} wprowadzany jest sygnał aktywności \texttt{<v>\_\_VforceEn<ID>}, który przechowuje informację o tym, na jakim przedziale bitowym dane przypisanie \texttt{force} jest aktualnie aktywne,
    \item w momencie wykonania instrukcji \texttt{force <v>[b:a]} odpowiedni sygnał \texttt{<v>\_\_VforceEn<ID>} jest ustawiany na~1 na przedziale \texttt{[b:a]},
    \item w momencie wykonania instrukcji \texttt{release <v>[b:a]} sygnały \texttt{<v>\_\_VforceEn<ID>} dla wszystkich identyfikatorów \texttt{ID} są zerowane na przedziale \texttt{[b:a]},
    \item każdy odczyt sygnału, do którego kiedykolwiek wykonywane było przypisanie \texttt{force}, jest zastępowany złożonym wyrażeniem iterującym po identyfikatorach \texttt{ID}, które na podstawie sygnałów \texttt{<v>\_\_VforceEn<ID>} oraz odpowiadających im wyrażeń z prawej strony konstruuje wartość wymuszoną.

\end{itemize}

\coderef{lst:conditional-force}~zostanie przekształcony na:
\begin{minted}[linenos]{systemverilog}
logic [1:0] a;
logic [1:0] a__VforceEn0 = 2'h0;
logic [1:0] a__VforceEn1 = 2'h0;
logic [1:0] b;
initial begin
    if (`IMPURE_ONE) begin
        a__VforceEn0 = 2'h3;
        a__VforceEn1 = a__VforceEn1 & (~ 2'h3);
    end

    if (~ `IMPURE_ONE) begin
        a__VforceEn1 = 2'h3;
        a__VforceEn0 = a__VforceEn0 & (~ 2'h3);
    end

    b = 2'h0;

    #1 $display("%d", ((a__VforceEn1 & b) | (a__VforceEn0 & 2'h1)) | ((~ (a__VforceEn1 | a__VforceEn0)) & a));

    a = ((a__VforceEn1 & b) | (a__VforceEn0 & 2'h1)) | ((~ (a__VforceEn1 | a__VforceEn0)) & a);
    a__VforceEn1 = a__VforceEn1 & 2'h0;
    a__VforceEn0 = a__VforceEn0 & 2'h0;
end
\end{minted}
\codecaption{lst:force-ast}{Przekształcenie wykonane przez Verilator z napisaną przeze mnie implementacją \texttt{force/release}, pozyskane przy użyciu flag \texttt{----dump-tree ----debug-emitv}, uproszczone i sformatyzowane ręcznie dla czytelności. Linie 7-8 odpowiadają pierwszemu przypisaniu \texttt{force}, linie 12-13 drugiemu. W linii 18 znajduje się złożone wyrażenie odczytujące sygnał, a linie 20-22 realizują instrukcję \texttt{release}.}

Ta implementacja poprawnie obsługuje przypadek warunkowego wykonywania instrukcji \texttt{force} oraz przechodzi z powodzeniem wszystkie testy z zestawu testów Verilatora, które używają instrukcji \texttt{force}/\texttt{release}.

Ma ona jednak bardzo poważną wadę: bardzo niską wydajność, zarówno czasową jak i pamięciową. Każde przypisanie \texttt{force} wymaga dodatkowej zmiennej \linebreak \texttt{<v>\_\_VforceEn<ID>}, a każde odczytanie sygnału z wymuszeniem wymaga złożonego wyrażenia. Każdy sygnał z przypisaniami \texttt{force} powoduje powiększenie drzewa AST o czynnik rzędu $O(n*(m+o))$ , gdzie $n$ jest liczbą przypisań \texttt{force}, $m$ liczbą odczytów sygnału, natomiast $o$ liczbą instrukcji \texttt{release} wykonywanych do tego sygnału gdy nie jest on siecią.

Wydawać by się mogło, że w praktyce zarówno liczba przypisań \texttt{force}/\texttt{release} do jednego sygnału jak i liczba jego odczytów nie powinny być duże. Tak rzeczywiście jest --- z tego powodu obszerny zestaw testów Verilatora nie sprawia tej implementacji większego problemu. Pojawia się jednak problem ze strukturami typu \texttt{struct packed} --- przed \texttt{V3Force} następuje przejście \texttt{V3Width}, które każdą taką strukturę zastępuje wektorem bitów. W przypadku struktury z wieloma polami oraz podstrukturami, taki wektor bitów może osiągać szerokość rzędu dziesiątek tysięcy bitów.

Przypadkiem projektu układu cyfrowego, który obnaża ten problem jest Caliptra,~\cite{caliptra} w której po transformacji dokonanej przez \texttt{V3Width} struktura \linebreak\texttt{caliptra\_top\_tb.caliptra\_top\_dut.key\_vault1.kv\_reg\_hwif\_in} jest zamieniana na wektor 13923 bitów, do którego jest wykonywane
\begin{itemize}
    \item 22586 przypisań \texttt{force},
    \item 2736 przypisań \texttt{release},
    \item 3630 odczytów.
\end{itemize}
W wyniku tego, próba verilacji Caliptry z tą implementacją potrafi skończyć błędem z powodu wykorzystania 512GB pamięci RAM w czasie poniżej 2 minut.

\section{Zastąpienie złożonych wyrażeń w drzewie pętlami}

Pomysłem na usprawnienie poprzedniej implementacji~\footnote{https://github.com/ArturBieniek4/verilator/tree/81646-force-rework-v4}, umożliwiając verilację i przetestowanie wydajności symulacji było zastąpienie złożonych wyrażeń w drzewie AST pętlami \texttt{while} (identycznymi jak te dostępne w Verilogu), które w czasie działania symulacji będą wykonywały odpowiednie przypisania i zwracały wartość wymuszoną. Ogranicza to rozrost drzewa AST na etapie verilacji kosztem zwiększenia liczby instrukcji wykonywanych w czasie symulacji.
Implementacja ta polega na stworzeniu:
\begin{itemize}
    \item tablicy wektorów bitowych \texttt{<v>\_\_VforceEns},
    \item tablicy sygnałów \texttt{wire} \texttt{<v>\_\_VforceRHS}.
\end{itemize}

Tablica \texttt{<v>\_\_VforceEns} przechowuje aktualny stan przypisań \texttt{force}, a tablica \texttt{<v>\_\_VforceRHS} przechowuje listę wszystkich prawych stron przypisań \texttt{force} do sygnału \texttt{<v>}.

Verilacja polega na:
\begin{itemize}
    \item stworzeniu globalnej zmiennej iteracyjnej \texttt{\_\_Vforcei} typu \texttt{integer signed} (32-bitowego) do wykorzystywania w pętlach \texttt{while} --- zmienna ta będzie zerowana przed każdą pętlą \texttt{while} i wykorzystywana tylko wewnątrz niej,
    \item zidentyfikowaniu wszystkich przypisań \texttt{force} do danego sygnału \texttt{<v>}, nadaniu im unikalnych identyfikatorów \texttt{ID} oraz utworzeniu struktur danych przechowujących wyrażenia z prawej strony tych przypisań, które miały odpowiednie przypisania kombinacyjne, w postaci tablicy \texttt{<v>\_\_VforceRHS} indeksowanej identyfikatorem \texttt{ID},
    \item zastąpieniu każdego przypisania \texttt{force <v>[b:a] = <RHS>} instrukcjami ustawiającymi sygnał aktywności przypisania o danym identyfikatorze \texttt{ID} na przedziale \texttt{[b:a]}, realizowanymi poprzez ustawienie odpowiednich bitów w tablicy \texttt{<v>\_\_VforceEns}, oraz jednoczesnym zerowaniem pokrywających się bitów odpowiadających pozostałym identyfikatorom, w celu zapewnienia, że na danym bicie aktywne jest co najwyżej jedno przypisanie \texttt{force},
    \item zastąpieniu każdego przypisania \texttt{release <v>[b:a]} instrukcjami zerującymi odpowiednie bity w tablicy \texttt{<v>\_\_VforceEns} dla wszystkich identyfikatorów \texttt{ID} na przedziale \texttt{[b:a]}, przy czym w przypadku, gdy \texttt{<v>} było siecią, wartość sygnału była wcześniej odczytywana,
    \item zastąpieniu każdego odczytu sygnału \texttt{<v>} złożonym wyrażeniem zawierającym pętlę \texttt{while}, która iterowała po identyfikatorach \texttt{ID} poprzez indeksację tablicy \texttt{<v>\_\_VforceEns}, konstruowała wartość wymuszoną jako alternatywę logiczną koniunkcji sygnałów aktywności i odpowiadających im wartości z tablicy \texttt{<v>\_\_VforceRHS}, budowała maskę bitów objętych wymuszeniem oraz zwracała wartość wynikową złożoną z alternatywy logicznej tej wartości i koniunkcji logicznej zanegowanej maski wymuszeń z wartością sygnału \texttt{<v>}.



\end{itemize}

\coderef{lst:conditional-force}~zostanie przez tę implementację przekształcony na:
\begin{minted}[linenos]{systemverilog}
logic [1:0] a;
logic [1:0] a__VforceRHS[1:0];
logic [1:0] a__VforceNeg;
logic [1:0] a__VforceRd;
logic [1:0] a__VforceEns[1:0] = '{default:'0};
logic [1:0] b;
integer signed __Vforcei;
initial a__VforceRHS[32'h0] = 2'h1;
initial a__VforceRHS[32'h1] = b;

always @(*) a__VforceRHS[32'h1] = b;
always @(*) begin
    a__VforceRd = 2'h0;
    a__VforceNeg = 2'h0;
    __Vforcei = 32'h0;
    while (__Vforcei < 32'h2) begin
        a__VforceRd = a__VforceRd | (a__VforceEns[__Vforcei] & a__VforceRHS[__Vforcei]);
        a__VforceNeg = a__VforceNeg | a__VforceEns[__Vforcei];
        __Vforcei = __Vforcei + 32'h1;
    end
    a__VforceRd = a__VforceRd | ((~a__VforceNeg) & a);
end

initial begin
    if (`IMPURE_ONE) begin
        a__VforceEns[32'h0] = 2'h3;
        a__VforceRHS[32'h0] = 2'h1;
        __Vforcei = 32'h0;
        while (__Vforcei < 32'h2) begin
            if (__Vforcei != 32'h0) begin
                a__VforceEns[__Vforcei] = a__VforceEns[__Vforcei] & (~ 2'h3);
            end
            __Vforcei = __Vforcei + 32'h1;
        end
    end

    if (~`IMPURE_ONE) begin
        a__VforceEns[32'h1] = 2'h3;
        a__VforceRHS[32'h1] = b;
        __Vforcei = 32'h0;
        while (__Vforcei < 32'h2) begin
            if (__Vforcei != 32'h1) begin
                a__VforceEns[__Vforcei] = a__VforceEns[__Vforcei] & (~ 2'h3);
            end
            __Vforcei = __Vforcei + 32'h1;
        end
    end

    b = 2'h0;

    #1;
    a__VforceRd = 0;
    a__VforceNeg = 0;
    __Vforcei = 0;
    while (__Vforcei < 2) begin
        a__VforceRd = a__VforceRd | (a__VforceEns[__Vforcei] & a__VforceRHS[__Vforcei]);
        a__VforceNeg = a__VforceNeg | a__VforceEns[__Vforcei];
        __Vforcei = __Vforcei + 1;
    end
    a__VforceRd = a__VforceRd | ((~a__VforceNeg) & a);
    $display("%d", a__VforceRd);

    a__VforceRd = 2'h0;
    a__VforceNeg = 2'h0;
    __Vforcei = 32'h0;
    while (__Vforcei < 32'h2) begin
        a__VforceRd = a__VforceRd | (a__VforceEns[__Vforcei] & a__VforceRHS[__Vforcei]);
        a__VforceNeg = a__VforceNeg | a__VforceEns[__Vforcei];
        __Vforcei = __Vforcei + 32'h1;
    end
    a__VforceRd = a__VforceRd | ((~a__VforceNeg) & a);
    a = a__VforceRd;
    __Vforcei = 32'h0;
    while (__Vforcei < 32'h2) begin
        a__VforceEns[__Vforcei] = 2'h0;
        __Vforcei = __Vforcei + 32'h1;
    end
    a__VforceRd = a;
end
\end{minted}
\codecaption{lst:force-loop}{Przekształcenie wykonane przez Verilator z napisaną przeze mnie implementacją obsługi \texttt{force}/\texttt{release}, pozyskane przy użyciu flag \texttt{----dump-tree ----debug-emitv}, uproszczone i sformatyzowane ręcznie dla czytelności. Linie 1-9 są odpowiedzialne za inicjalizację tablicy \texttt{<v>\_\_VforceRHS} oraz zerowanie tablicy \texttt{<v>\_\_VforceEns}. Linie 11-22 odpowiadają za ciągłe przypisanie wartości sygnału \texttt{<v>\_\_VForceRd}, analogicznie jak we wcześniej omawianej obecnej implementacji. Linie 26-34 realizują pierwsze przypisanie \texttt{force}, linie 38-46 drugie przypisanie \texttt{force}. W liniach 52-60 znajduje się aktualizacja wartości sygnału \texttt{<v>\_\_VforceRd} przed wyświetleniem jego wartości, a linie 63-78 realizują instrukcję \texttt{release}.}

Ta implementacja również poprawnie obsługuje przypadek warunkowego wykonywania instrukcji \texttt{force} oraz przechodzi z powodzeniem wszystkie testy z zestawu testów Verilatora, które używają instrukcji \texttt{force}/\texttt{release}.

Złożoność pamięciowa zaproponowanej implementacji w fazie verilacji jest liniowa względem łącznego rozmiaru struktur danych tworzonych dla obsługi instrukcji \texttt{force}. Dla każdego sygnału $v$ objętego przypisaniami \texttt{force} tworzona jest tablica \texttt{<v>\_\_VforceEns} oraz tablica \texttt{<v>\_\_VforceRHS}, których rozmiary są proporcjonalne do iloczynu liczby przypisań \texttt{force} dla danego sygnału, oznaczonej jako $F_v$, oraz jego szerokości bitowej $W_v$. W konsekwencji całkowite zużycie pamięci w czasie verilacji można opisać jako
\[
    \Theta\!\left(\sum_v F_v \cdot W_v\right),
\]
co odpowiada liniowej złożoności względem łącznej liczby bitów koniecznych do reprezentacji wszystkich przypisań \texttt{force} w projekcie.

Złożoność czasowa w fazie verilacji jest również liniowa względem rozmiaru danych wejściowych. Każde wystąpienie instrukcji \texttt{force}, \texttt{release} oraz każdy odczyt sygnału są analizowane i transformowane dokładnie jednokrotnie, a koszt pojedynczej transformacji jest proporcjonalny do rozmiaru przetwarzanego fragmentu, w szczególności do szerokości sygnału oraz rozmiaru prawej strony przypisania. W rezultacie całkowity czas verilacji jest proporcjonalny do sumy liczby instrukcji \texttt{force}/\texttt{release}, liczby odczytów sygnałów oraz łącznego rozmiaru przetwarzanych wyrażeń, co potwierdza liniową złożoność czasową zaproponowanego rozwiązania.

Problemem tej implementacji pozostaje jednak złożoność czasowa w czasie działania symulacji. W najgorszym przypadku, każde odczytanie sygnału z wymuszeniem wymaga iteracji po wszystkich przypisaniach \texttt{force} do tego sygnału, co prowadzi do złożoności czasowej $O(m \cdot n)$, gdzie $m$ jest liczbą odczytów sygnału, a $n$ liczbą przypisań \texttt{force} do tego sygnału. W praktyce, dla większości testów z zestawu testów Verilatora, liczba przypisań \texttt{force} do jednego sygnału jest niewielka, więc wydajność symulacji pozostaje akceptowalna. Jednak w przypadku projektów takich jak Caliptra, gdzie liczba przypisań \texttt{force} do jednego sygnału może być znaczna, wydajność symulacji ulega drastycznemu pogorszeniu. W przypadku Caliptry, próba symulacji z tą implementacją skutkuje czasem symulacji przekraczającym 3 godziny, co jest nieakceptowalne w praktycznych zastosowaniach.

\section{Implementacja za pomocą struktur danych w C++}

Inną implementacją\footnote{https://github.com/ArturBieniek4/verilator/tree/89452-force-rework-v5}, którą opracowałem, jest podejście oparte na strukturach danych zaimplementowanych w C++ i wykorzystywanych w czasie działania symulacji. W tej implementacji wykorzystałem struktury do reprezentacji stanów sygnałów oraz ich wymuszeń, co pozwoliło na uproszczenie logiki przetwarzania i zwiększenie wydajności. Zaimplementowana struktura wykorzystuje posortowany wektor do przechowywania aktywnych przypisań \texttt{force}, co umożliwia szybkie wyszukiwanie i aktualizację przypisań. Każde przypisanie \texttt{force} jest reprezentowane jako obiekt zawierający informacje o zakresie bitów, które są objęte tym przypisaniem oraz wartości która jest przypisana. Taki wektor jest trzymany zawsze jako posortowany według początków zakresów bitów, co umożliwia wyszukiwanie pasujących przypisań za pomocą wyszukiwania binarnego w czasie $O(log\ k)$, gdzie $k$ jest liczbą aktywnych przypisań \texttt{force}.
W momencie wykonania instrukcji \texttt{release}:
\begin{itemize}
    \item jeżeli lewa strona instrukcji nie jest siecią, to dokonywany jest odczyt jej aktualnej wartości, która zostaje przypisana do sygnału,
    \item znajdowane są pasujące przypisania \texttt{force} w wektorze aktywnych przypisań,
    \item granice pasujących przypisań są modyfikowane, a jeśli któreś z nich jest całkowicie zwolnione, to jest usuwane z wektora.
\end{itemize}

Takie podejście zapewnia liniową złożoność czasową i pamięciową w czasie verilacji, zależną jedynie od sumy liczby przypisań \texttt{force} i \texttt{release}, wartości po ich prawych stronach oraz liczby odczytów sygnału --- czyli taką samą jak w poprzedniej implementacji.

Złożoność pamięciowa w czasie działania symulacji jest również liniowa, względem sumy liczby aktywnych przypisań \texttt{force} oraz rozmiarów ich prawych stron, co jest znacznie bardziej efektywne niż poprzednia implementacja.

Złożoność czasowa pojedynczego odczytu sygnału, przypisania \texttt{force} do sygnału lub instrukcji \texttt{release} w czasie symulacji wynosi $O(log\ k)$, gdzie $k$ jest liczbą aktywnych przypisań \texttt{force} do danego sygnału. Wynika to z faktu, że wyszukiwanie pasujących przypisań w posortowanym wektorze odbywa się za pomocą wyszukiwania binarnego, a dodatkowe operacje wykonywane po znalezieniu pasujących przypisań mają stały koszt czasowy.

\coderef{lst:conditional-force}~zostanie przekształcony na:
\begin{minted}[linenos]{systemverilog}
logic [1:0] a;
logic [1:0] a__VforceRHS[1:0];
VlForceVec a__VforceVec;
logic [1:0] b;
always @(*) a__VforceRHS[32'h1] = b;
always @(*) a__VforceRHS[32'h0] = 2'h1;

initial begin
    if (`IMPURE_ONE) begin
        a__VforceRHS[32'h0] = 2'h1;
        a__VforceVec.addForce(0, 1, 0);
    end

    if (~`IMPURE_ONE) begin
        a__VforceRHS[32'h1] = b;
        a__VforceVec.addForce(0, 1, 1);
    end

    b = 2'h0;

    #1 $display("%d", VlForceRead(a, a__VforceVec, a__VforceRHS));

    a = VlForceRead(a, a__VforceVec, a__VforceRHS);
    a__VforceVec.release(0, 1);
end
\end{minted}
\codecaption{lst:force-cpp}{Przekształcenie wykonane przez Verilator z napisaną przeze mnie implementacją obsługi \texttt{force}/\texttt{release}, pozyskane przy użyciu flag \texttt{----dump-tree ----debug-emitv}, uproszczone i sformatyzowane ręcznie dla czytelności. Linie 1-6 są odpowiedzialne za inicjalizację tablicy \texttt{<v>\_\_VforceRHS} oraz struktury \texttt{<v>\_\_VforceVec}. Linie 10-11 odpowiadają pierwszemu przypisaniu \texttt{force}, linie 15-16 drugiemu. Linia 21 odczytuje sygnał, a linie 23-24 realizują instrukcję \texttt{release}.}

Jak widać w powyższym fragmencie kodu, instrukcje \texttt{force} są zamieniane na wywołania metody \texttt{addForce} klasy \texttt{VlForceVec}, natomiast instrukcje \texttt{release} są zamieniane na wywołania metody \texttt{release} tej klasy. Odczyty sygnałów do których wykonywane są przypisania \texttt{force} są zamieniane na wywołania funkcji \texttt{VlForceRead}, która przyjmuje oryginalną wartość sygnału, obiekt klasy \texttt{VlForceVec} przechowujący informacje o aktywnych przypisaniach \texttt{force} oraz wskaźnik na tablicę prawych stron przypisań \texttt{force}.

Poniżej znajduje się porównanie wydajności wykonane przy użyciu narzędzia RTLMeter~\cite{rtlmeter}. Najnowszy dostępny Verilator jest porównywany z wyżej opisaną implementacją na niego nałożoną.

Pomiary zostały wykonane na komputerze z procesorem AMD Ryzen 9 7900X oraz 32GB pamięci RAM DDR5-6000MT/s. SMT (\textit{Simultaneous Multithreading})  oraz boost były wyłączone, a governor częstotliwości procesora ustawiony na \textit{performance}. Każdy z testów został uruchomiony 10 razy, a podane wyniki to mediana z tych uruchomień.

\begin{table}[H]
\caption{Porównanie wydajności pełnej verilacji (verilate + cppbuild + compile) przed i po wprowadzeniu zmian}
\small
\begin{tabular*}{\linewidth}{@{\extracolsep{\fill}}|l|c|c|c|c|c|c|}
\hline
& \multicolumn{2}{c|}{Referencja}
& \multicolumn{2}{c|}{Implementacja}
& \multicolumn{2}{c|}{Zmiana [\%]} \\
\hline
Design
& Czas [s] & RAM [MB]
& Czas [s] & RAM [MB]
& Czas [\%] & RAM [\%] \\
\hline
\benchrow{NVDLA\cite{NVDLA}}{1429.61}{8114.85}{1427.39}{8115.90}
\benchrow{OpenPiton\cite{OpenPiton}}{504.08}{2699.64}{504.32}{2699.69}
\benchrow{OpenTitan\cite{opentitan}}{270.86}{1784.65}{271.00}{1785.01}
\benchrow{VeeR\cite{veer_eh2}}{143.14}{933.43}{144.21}{933.40}
\benchrow{Vortex\cite{vortex}}{239.84}{1075.07}{237.48}{1075.66}
\benchrow{XuanTie\cite{XUANTIE_C910}}{424.15}{2542.06}{423.44}{2542.33}
\end{tabular*}
\label{tab:verilator-compare-time-mem}
\end{table}


\begin{table}[H]
\caption{Porównanie wydajności symulacji przed i po wprowadzeniu zmian}
\small
\begin{tabular*}{\linewidth}{@{\extracolsep{\fill}}|l|c|c|c|c|c|c|}
\hline
& \multicolumn{2}{c|}{Referencja}
& \multicolumn{2}{c|}{Implementacja}
& \multicolumn{2}{c|}{Zmiana [\%]} \\
\hline
Design
& Czas [s] & RAM [MB]
& Czas [s] & RAM [MB]
& Czas [\%] & RAM [\%] \\
\hline
\benchrow{NVDLA}{222.26}{1238.69}{222.03}{1238.65}
\benchrow{OpenPiton}{64.41}{41.77}{64.40}{41.74}
\benchrow{OpenTitan}{141.71}{22.62}{142.45}{22.51}
\benchrow{VeeR}{54.97}{13.56}{54.12}{13.64}
\benchrow{Vortex}{38.41}{19.65}{38.37}{19.67}
\benchrow{XuanTie}{88.95}{61.33}{89.11}{61.36}

\end{tabular*}
\label{tab:verilator-compare-time-mem-execute}
\end{table}

Wyniki pokazują, że wprowadzone zmiany \textbf{nie wpływają istotnie na zużycie pamięci ani czas verilacji oraz symulacji} — różnice są pomijalne i mieszczą się w granicach błędu pomiarowego.

\chapter{Dalsza praca}

\section{Wsparcie dla \textit{assign/deassign}}
Z poprawnie działającą implementacją instrukcji \texttt{force/release}, kolejnym krokiem jest implementacja instrukcji \texttt{assign/deassign}. Jak wspomniano wcześniej, instrukcje te realizują identyczne funkcje jak \texttt{force/release}, ale zbiór sygnałów które mogą być ich lewą stroną jest podzbiorem sygnałów które mogą być lewą stroną \texttt{force/release}. W szczególności, instrukcje \texttt{assign/deassign} nie mogą być użyte do wymuszenia wartości na sieciach oraz nie mogą posiadać złożonych lewych stron przypisań.

Z tego powodu, implementacja instrukcji \texttt{assign/deassign} może być zrealizowana jako specjalny, ograniczony przypadek implementacji instrukcji \texttt{force/release}. To podejście zostało już przetestowane z pierwszą opisaną implementacją instrukcji \texttt{force/release} i okazało się skuteczne. Po złączeniu docelowej implementacji \texttt{force}/\texttt{release} do głównej gałęzi Verilatora, kolejnym krokiem będzie implementacja tego podejścia.

\section{Wsparcie dla \texttt{force/release}, gdzie prawą stroną jest wywołanie funkcji}

Kolejnym krokiem jest dodanie wsparcia dla przypadków, gdy prawa strona instrukcji \texttt{force} zawiera wywołania funkcji czystych\footnote{Funkcja czysta (ang. \textit{pure}) to taka funkcja, która nie ma efektów ubocznych ani wewnętrznego stanu i zawsze zwraca tę samą wartość dla tych samych argumentów. Dokładna definicja funkcji czystej w SystemVerilogu znajduje się w sekcji 35.5.2 IEEE 1800-2023~\cite{IEEE1800-2023}.}. W obecnej implementacji, jeśli prawa strona instrukcji \texttt{force} jest wywołaniem funkcji, to Verilator tworzy zmienną, której przypisywana jest wartość zwrócona przez funkcję w momencie wykonania instrukcji \texttt{force}, a następnie ta zmienna jest używana jako prawa strona przypisania. To podejście jest oczywiście niepoprawne, ponieważ wartość zwrócona przez funkcję może się zmieniać w czasie symulacji, podczas gdy przypisanie \texttt{force} powinno zawsze używać aktualnej wartości zwróconej przez funkcję w momencie odczytu sygnału.

Prostym rozwiązaniem tego problemu jest stworzenie dodatkowego sygnału typu \texttt{wire} dla każdego wywołania funkcji użytego jako prawa strona instrukcji \texttt{force}, który będzie miał aktywne przypisanie kombinacyjne do wartości zwróconej przez funkcję. Następnie, w momencie wykonania instrukcji \texttt{force}, ten sygnał będzie używany jako prawa strona przypisania. W ten sposób, przy każdym odczycie sygnału z aktywnym przypisaniem \texttt{force}, aktualnie zwracana wartość będzie brana pod uwagę.

Oczywiście użycie z prawej strony instrukcji \texttt{force} funkcji nieczystych powinno skutkować błędem. To zachowanie jest już zaimplementowane w Verilatorze i użycie przypisania innego niż proceduralne z wywołaniem funkcji nieczystej skutkuje błędem kompilacji.

To podejście również zostało już przetestowane z pierwszą opisaną implementacją instrukcji \texttt{force/release} i okazało się skuteczne. Po złączeniu docelowej implementacji \texttt{force}/\texttt{release} do głównej gałęzi Verilatora, kolejnym krokiem będzie implementacja tego podejścia.

\section{Rozszerzenie interfejsu DPI-C}
Verilator wspiera interfejs DPI-C, umożliwiający sterowanie symulacją z poziomu kodu C/C++. Interfejs ten ma możliwość sterowania przypisaniami \texttt{force/release} poprzez udostępnienie sygnałów \texttt{<v>\_VforceEn} i \texttt{<v>\_VforceVal} dla każdego sygnału \texttt{<v>}, do którego jest wykonywane przypisanie \texttt{force}.
    
Poprzednia implementacja instrukcji \texttt{force/release} w Verilatorze wykorzystywała te sygnały do realizacji wymuszania wartości, co dawało również możliwość odczytu i modyfikacji stanu przypisania \texttt{force} z poziomu kodu C/C++.

\begin{minted}{systemverilog}
module top;
    wire [1:0] out /* verilator forceable */;
    initial begin
        force out = 0;
        $display(out);
    end
endmodule
\end{minted}
\codecaption{lst:dpi-force}{Przykład kodu, który pozwala na sterowanie przypisaniami \texttt{force} z poziomu kodu C/C++ poprzez interfejs DPI-C.}

W docelowej implementacji, sygnały te nie są już używane do realizacji wymuszania wartości, jednak zostały utrzymane w celu zachowania kompatybilności z istniejącym kodem korzystającym z interfejsu DPI-C.
Sygnały te mają obecnie wyższy priorytet niż przypisania \texttt{force} wykonywane w kodzie (System)Verilog. Realizowane są poprzez dodanie sygnałów takich jak w pierwotnej implementacji, z identyczną mechaniką działania.

\begin{minted}{systemverilog}
logic [1:0] out;
logic [1:0] out__VforceRHS[1:0];
VlForceVec out__VforceVec;
logic [1:0] out__VforceVal;
logic [1:0] out__VforceEn;
always @(out__VforceEn) begin
    out__VforceRHS[32'h0] = (out__VforceEn & out__VforceVal) | ((~ out__VforceEn) & out);
    out__VforceVec.addForce(0, 1, 0);
end
always @(*) out__VforceRHS[32'h1] = 2'h0;
always @(*) out__VforceRHS[32'h0] = (out__VforceEn & out__VforceVal) | ((~ out__VforceEn) & out);
initial begin
    out__VforceRHS[32'h1] = 2'h0;
    out__VforceVal = (out__VforceVal & (~ 2'h3)) | 2'h0;
    out__VforceEn = out__VforceEn | 2'h3;
    out__VforceVec.addForce(0, 1, 1);
    $display("%d", VlForceRead(out, out__VforceVec, out__VforceRHS));
end
\end{minted}
\codecaption{lst:dpi-force-transformation}{Przekształcenie wykonane przez Verilator, pozyskane przy użyciu flag \texttt{----dump-tree ----debug-emitv}, uproszczone i sformatyzowane ręcznie dla czytelności.}

Indeks zerowy w tablicy \texttt{<v>\_\_VforceRHS} jest wówczas zawsze rezerwowany dla DPI-C. Tworzone są publicznie dostępne zmienne \texttt{<v>\_\_VforceVal} oraz \texttt{<v>\_\_VforceEn}. Zmiana \texttt{<v>\_\_VforceEn} powoduje aktualizację \texttt{<v>\_\_VforceRHS[32'h0]} oraz ponowne włączenie tego wymuszenia za pomocą funkcji \texttt{addForce}. Takie podejście pozwala zachować kompatybilność z istniejącymi testami, nie wpływając znacząco na wydajność symulacji.

Pomysłem na dalsze rozszerzenie interfejsu DPI-C jest udostępnienie dodatkowo sygnału \texttt{<v>\_VforceVec} i/lub metod \texttt{addForce} oraz \texttt{release}, co da dodatkowo możliwość zarządzania wszystkimi przypisaniami \texttt{force} wykonywanymi z poziomu kodu (System)Verilog.

Ta funkcjonalność może być zaimplementowana po złączeniu docelowej implementacji \texttt{force}/\texttt{release} do głównej gałęzi Verilatora. To, czy podjęta będzie próba jej implementacji, zależy od zainteresowania społeczności Verilatora tym rozszerzeniem i poziomu trudności tej implementacji.

\chapter{Podsumowanie i wnioski}
\section{Najważniejsze rezultaty}

W ramach niniejszej pracy przeprowadzono szczegółową analizę semantyki instrukcji \texttt{force}/\texttt{release} w odniesieniu do standardu IEEE, ze szczególnym uwzględnieniem jej relacji do innych mechanizmów przypisań występujących w językach Verilog i SystemVerilog, takich jak przypisania kombinacyjne, proceduralne oraz proceduralne kombinacyjne. Analiza ta pozwoliła na jednoznaczne określenie oczekiwanego, obserwowalnego zachowania sygnałów w obecności wymuszeń oraz na identyfikację przypadków problematycznych z punktu widzenia implementacji w symulatorach opartych na kompilacji.

Następnie przeanalizowano istniejący w Verilatorze mechanizm obsługi instrukcji \texttt{force}/\texttt{release}, zaimplementowany w pliku \texttt{V3Force.cpp}. Wskazano zarówno jego zalety, wynikające z prostoty zastosowanej transformacji i korzystnego wpływu na wydajność, jak i istotne ograniczenia funkcjonalne. W szczególności wykazano problemy z poprawnością semantyczną w przypadku warunkowego wykonywania instrukcji \texttt{force}, wynikające z przyjętych założeń statycznej interpretacji.

W odpowiedzi na zidentyfikowane problemy zaprojektowano i porównano kilka alternatywnych podejść implementacyjnych, których celem było zapewnienie pełnej zgodności z semantyką języka. Zaproponowane rozwiązania obejmowały m.in. podejście polegające na zastępowaniu odczytów sygnałów złożonymi wyrażeniami zależnymi od aktualnego stanu wymuszeń, jak również wariant redukujący rozrost drzewa AST poprzez przeniesienie części logiki do pętli wykonywanych dynamicznie w trakcie symulacji. Ostatecznie opracowano docelową implementację wykorzystującą dedykowane struktury danych w języku C++ do reprezentacji aktywnych wymuszeń oraz wyspecjalizowane funkcje odczytu, które łączą wartość wynikającą z normalnej logiki projektu z nałożonymi wymuszeniami.

Poprawność zaprojektowanych rozwiązań została zweryfikowana przy użyciu testów regresyjnych Verilatora, a także poprzez analizę zachowania w wybranych przypadkach brzegowych, w tym w projektach intensywnie wykorzystujących złożone struktury danych, takie jak \texttt{packed struct}. Pozwoliło to potwierdzić zgodność implementacji z oczekiwaną semantyką języka również w mniej typowych scenariuszach użycia.

Na końcu oceniono wpływ docelowego podejścia na wydajność procesu verilacji oraz symulacji, przeprowadzając pomiary na reprezentatywnych projektach z wykorzystaniem narzędzia RTLMeter. Uzyskane wyniki wskazują na brak istotnego wzrostu zużycia pamięci oraz na jedynie niewielkie, zależne od charakterystyki projektu różnice czasowe, co potwierdzają zestawienia przedstawione w Tab.~\ref{tab:verilator-compare-time-mem} oraz Tab.~\ref{tab:verilator-compare-time-mem-execute}.

Dodatkowo, w toku prac wprowadzono poprawki zwiększające jakość i zakres obsługiwanych przypadków w istniejącym mechanizmie, w tym m.in. rozszerzenia dotyczące wyrażeń z prawej strony \texttt{force}, diagnostyki błędów oraz korekty związanej z dopasowaniem szerokości operandów.

\section{Wnioski z porównania podejść implementacyjnych}

Przeprowadzone eksperymenty pokazują, że w kontekście Verilatora poprawna obsługa \texttt{force}/\texttt{release} wymaga utrzymywania w czasie symulacji informacji o \emph{dynamicznym} stanie wymuszeń, a nie jedynie statycznego „ostatniego \texttt{force}” znanego na etapie kompilacji. Najważniejsze obserwacje są następujące:

\begin{itemize}
    \item \textbf{Rozwiązania oparte wyłącznie o transformacje AST mogą być poprawne, ale nie muszą być skalowalne.}
    Implementacja oparta o zastępowanie odczytów złożonymi wyrażeniami jest koncepcyjnie prosta i poprawna semantycznie, jednak prowadzi do gwałtownego rozrostu drzewa AST w projektach, w których (po transformacjach) pojedyncze sygnały osiągają bardzo dużą szerokość, a \texttt{force}/\texttt{release} występują masowo. Taki scenariusz pojawia się w praktyce m.in. przy intensywnym użyciu struktur typu \texttt{struct packed}.
    
    \item \textbf{Przeniesienie części pracy do czasu symulacji ogranicza zużycie pamięci na etapie kompilacji, ale może być nieakceptowalne czasowo.}
    Zastąpienie rozbudowanych wyrażeń pętlami wykonywanymi podczas symulacji poprawia złożoność pamięciową i umożliwia przetworzenie przypadków, które wcześniej były niepraktyczne na etapie verilacji. Jednak koszt wykonywania pętli przy każdym odczycie sygnału z wymuszeniem może dominować czas symulacji, szczególnie gdy liczba aktywnych lub potencjalnych wymuszeń jest bardzo duża.
    
    \item \textbf{Najbardziej zrównoważone jest podejście hybrydowe: minimalna ingerencja w kod wynikowy + struktury danych używane w symulacji.}
    Implementacja wykorzystująca struktury danych w C++ (np. uporządkowane wektory opisujące aktywne zakresy wymuszeń) pozwala ograniczyć rozrost kodu generowanego na etapie kompilacji, a jednocześnie zapewnia efektywny dostęp do informacji o wymuszeniach w czasie symulacji. W szczególności umożliwia to realizację operacji takich jak dodanie wymuszenia, zwolnienie zakresu oraz odczyt wartości z narzutem rzędu $O(\log k)$ względem liczby aktywnych wpisów, zamiast kosztów liniowych zależnych od liczby \texttt{force}.
\end{itemize}

Otrzymane wyniki wydajnościowe wskazują, że zaproponowana docelowa implementacja może być wdrożona bez istotnej regresji w typowych zastosowaniach: różnice w zużyciu pamięci są pomijalne, a wpływ na czas verilacji i symulacji jest niewielki oraz zależny od charakterystyki projektu (Tab.~\ref{tab:verilator-compare-time-mem} i Tab.~\ref{tab:verilator-compare-time-mem-execute}). Jest to istotne z punktu widzenia użytkowników Verilatora, ponieważ \texttt{force}/\texttt{release} są często używane w środowiskach weryfikacyjnych, gdzie oczekuje się zarówno poprawności semantycznej, jak i wysokiej wydajności.

\section{Podsumowanie}

W pracy wykazano, że poprawna i jednocześnie praktyczna wydajnościowo implementacja \texttt{force}/\texttt{release} w Verilatorze jest możliwa, mimo że narzędzie nie realizuje klasycznego modelu zdarzeniowego. Kluczowe okazało się odejście od wyłącznie statycznej interpretacji wymuszeń na etapie kompilacji i wprowadzenie mechanizmu, który utrzymuje stan wymuszeń w czasie symulacji oraz umożliwia efektywny odczyt wartości uwzględniający priorytety przypisań.

Ostatecznie zaproponowane podejście hybrydowe (transformacja do wywołań funkcji podczas symulacji + struktury danych w C++) zapewnia poprawność semantyczną w przypadkach, w których wcześniejsze rozwiązania były błędne, a równocześnie unika problemów skalowalności typowych dla podejść powodujących rozrost AST lub wymagających kosztownych pętli przy każdym odczycie. Wyniki pomiarów wskazują, że uzyskana funkcjonalność może zostać dostarczona użytkownikom bez istotnego pogorszenia parametrów praktycznych procesu verilacji i symulacji, co stanowi ważny krok w kierunku pełniejszego wsparcia konstrukcji SystemVerilog w Verilatorze.


%%%%% BIBLIOGRAFIA

\printbibliography

\end{document}
